% =============================================================================
% §8 Conclusion — arc restated; ARG/CHL forward hook; light policy resonance;
% ENDS as the resolving next step.
% =============================================================================
\section{Conclusion}
\label{sec:conclusion}

Latin America's fertility collapse is among the fastest on record, and fast
phenomena attract fast-dynamics explanations. We took the leading one
seriously enough to build it: a coordination-threshold model of reflexive norm
cascade, calibrated to the union-composition data. The data rejected it twice
--- generatively (a glide where the observations accelerate) and
econometrically (state dependence with the stabilizing sign, in every
specification, in both countries, at cell and individual-outcome level). What
survived is almost embarrassingly simple: successive cohorts arrive with
durably lower marriage propensity, and the observed collapse is the stock
composition catching up with the flow --- a compressed Second Demographic
Transition running through a consensual-union channel that was already wide.
A cohort microsimulation whose entire mechanism is four interpretable
parameters reproduces the collapse's magnitude and late-loading, and its one
recovered inflection point independently matches the econometric cohort
crossover. The exciting mechanism failed; the boring one, tested to
pre-registered standards, suffices.

Three things follow. For measurement: the collapse is invisible in smoothed
international series, and analysis or policy built on them will date it late
and understate it --- national registries and microdata are not optional here.
For interpretation: composition is the state variable. Where marriage and
cohabitation carry different fertility, a marriage collapse at flat union
totals is a fertility event, and projections that model union formation
without its composition will miss the mechanism. For policy, stated at the
modest level the evidence supports: the driver is cohort replacement, not a
within-cohort contagion --- interventions premised on interrupting a cascade
target a mechanism the data do not exhibit, while the actual structure implies
persistence, because cohorts already formed carry their propensities forward.

Two tests will discipline what we have claimed. The Argentine and Chilean
composition series will test the channel prediction out of sample: if the
consensual-union channel is the compression mechanism, marriage-dominant
regimes should collapse differently. And event-history data (the Colombian
ENDS program) will observe the transitions that repeated cross-sections cannot
--- resolving the one unreproduced shape residual, testing the excluded 15--19
entry margin, and providing the only evidence that could legitimately reopen
the reflexivity question. We have registered the predictions; the data will
adjudicate.
